\chapter{Union-Find}
\label{chap:union-find}

% Closed question: How do multiple records combine without collapsing into one
% undifferentiated record?

\begin{chapteressay}

Two physicists run an experiment together. Each makes notes in their
own lab notebook. At the end of the day, they sit down to consolidate
their records. The consolidation is not the trivial concatenation of
two notebooks. Each physicist has observed the same events from a
different vantage; each has notes about the same equipment, the same
samples, the same data, in different formats. The consolidation
requires merging records about the same event from two sources while
keeping each source identifiable.

If the consolidation is done badly, the result is a heap: a pile of
notes from which neither physicist can recover their original
observations. Specific facts become impossible to attribute. Errors
become impossible to localize. The consolidated notebook is less
useful than either original.

If the consolidation is done well, the result is a structured record
in which each event is identified across the two notebooks. ``Sample
A, temperature reading at 14:32, observed by both physicists. Notebook
1 records 22.3 C; Notebook 2 records 22.4 C; the discrepancy of 0.1 C
is within the thermometer's calibrated precision; the consolidated
value is 22.35 C with uncertainty 0.05 C.'' The consolidated record
preserves the source identity of each observation while presenting the
unified picture.

This is the chapter's question in its working form. Multiple records
must combine into a joint structure without losing their source
identities. The data structure that supports this is the union-find.

Union-find is the canonical algorithmic form for combining records
while preserving identity. The data structure maintains, for each
record, a pointer to a representative of its component. Records in
the same component are connected (in some specified sense); records
in different components are not. The structure supports two operations:
\texttt{union}, which merges two components into one, and \texttt{find},
which returns the representative of a record's component.

The chapter develops five instances of this structure. The first is
network connectivity: given a set of servers and connections, is
server A reachable from server Z? The union-find data structure
answers this in nearly constant time per query, after a one-time
$O(E + V)$ setup. The second is path compression: a small optimization
that flattens the union-find tree, reducing the cost of \texttt{find}
to amortized $O(\alpha(n))$ --- effectively constant for all practical
$n$. The third is the compiler symbol table: namespace operations in
Lean perform union-find-like bookkeeping of aliases and scope
resolution that preserve the source identity of definitions across
qualifying contexts. (The analogy is structural, not literal: Lean's
namespace machinery uses tries, scoping rules, and import graphs
rather than a literal disjoint-set forest, but the discipline ---
merge aliases without collapsing identity --- is the same.) The
fourth is the horizon as
union failure: causal boundaries that prevent merging, formalized in
the algorithm as components that the union operation refuses to
combine. The fifth is the typeclass machinery \texttt{Jar} /
\texttt{MeesaProcess} / \texttt{GUNGAN} in Lean: the formal
representation of unresolved component membership and its admissible
handling.

The second concrete example is the Regional Transit Authority. Three
predecessor agencies (city bus, light rail, inter-regional express)
merge into one unified system. Each predecessor agency has its own
fare cards and account histories. The unified system must combine
the three into a single passenger record without dissolving the
predecessor identities. The technical solution is a mapping table:
each unified pass has an associated set of predecessor cards. The
mapping is a union-find. The unified pass is the representative.
The predecessor cards are the members. Historical queries route
through the mapping to the correct predecessor database. The
predecessor identities are preserved; the unified pass is the
component.

If the merger had been done badly --- if the predecessor records had
been dissolved into anonymous unified records --- the audit trail
would be lost. A passenger asking ``what was my city bus balance
six months ago?'' would have no answer; the city bus balance no
longer exists as a distinct quantity. With proper union-find, the
question is answerable: the system traces the unified pass back to
the predecessor city bus card and returns the recorded balance.

The same structure appears in many domains. A compiler's symbol
table maps multiple names (qualified, unqualified, aliased) to the
same definition. A database's foreign-key system maps many records
in dependent tables to a single record in the primary table. A
version control system maps many file states across branches to a
single object hash. In each case, the algorithm is union-find. The
specific operations and representations vary; the structural
pattern is the same.

The chapter's closed question is:

\begin{center}
\emph{How do multiple records combine without collapsing into one
undifferentiated record?}
\end{center}

The answer the chapter develops is the union-find data structure
with explicit handling of component identity. The merge operation
combines records into components. The find operation retrieves the
component representative. Path compression keeps the queries fast.
Horizon failures are admitted as records of inadmissible mergers.
Unresolved component membership is handled through the
\texttt{Jar} / \texttt{MeesaProcess} / \texttt{GUNGAN} cascade.

The chapter's broader claim is that this is the algorithmic form
forced by the question. Any system that combines records while
preserving identities reduces to union-find. The data structure may
be implemented differently (as pointer trees, as hash tables, as
flat arrays with indices), but the algorithmic content is the same.
Combining without collapsing is union-find. There is no other
admissible algorithmic primitive for the task.

\end{chapteressay}

\section{Network Connectivity}
\label{se:network-connectivity}

A data center has ten servers and fifteen network connections. Each
connection links two servers. The question is: is server A reachable
from server Z, by any path through the available connections? The naive
approach is to perform a breadth-first or depth-first search starting at
A and see if Z is visited. The cost is $O(V + E)$ per query, where $V$
is the number of servers and $E$ is the number of connections. For ten
servers and fifteen connections, this is cheap. For ten million servers
and one hundred million connections, this is expensive.

The standard algorithm is union-find. The data structure maintains, for
each server, a pointer to a representative server in the same connected
component. Two servers are in the same component if and only if their
representatives are the same. The data structure supports two operations:
\texttt{union(a, b)}, which merges the components containing $a$ and
$b$, and \texttt{find(a)}, which returns the representative of $a$'s
component.

The initial state is that each server is its own component. Each server
points to itself. There are ten components.

When the connections are added, the algorithm calls \texttt{union(a, b)}
for each connection $(a, b)$. The union operation finds the
representatives of $a$ and $b$ (call them $r_a$ and $r_b$). If they are
the same, $a$ and $b$ are already in the same component; nothing changes.
If they are different, the algorithm picks one --- say $r_a$ --- and
updates $r_b$'s pointer to point to $r_a$. The two components are now
merged.

After processing all fifteen connections, the data structure represents
the connectivity of the network. Some servers may be in the same
component; some may not. To answer the connectivity question for $A$ and
$Z$, the algorithm computes \texttt{find(A)} and \texttt{find(Z)}. If
the results are equal, $A$ and $Z$ are connected; if not, they are not.

The cost of \texttt{find} depends on the depth of the pointer chain.
In the worst case, with naive union, the chain can be $O(V)$ deep, and
the total cost over many operations is $O(V^2)$. With union-by-rank
(always attaching the shorter tree to the root of the taller), the depth
is bounded by $O(\log V)$. With path compression (which the next section
develops), the amortized cost per operation is nearly $O(1)$.

The interesting property of union-find is its preservation of identity.
Two servers $A$ and $B$ that are merged into the same component are not
identified with each other; they remain distinct servers. The merging
identifies the components, not the elements. \texttt{find(A)} returns
the representative of $A$'s component, not $A$ itself. The component is
the equivalence class; the elements are the members. The component has a
representative; the elements have identities.

This is the chapter's question in its purest form. Multiple records
(servers) combine into a joint structure (the connected component)
without losing their individual identities. The records are connected;
they are not collapsed. The connectivity is one fact about them; the
remaining facts (each server's name, its hardware, its workload) are
preserved.

The Lean analog in \texttt{Episode4.lean} is the
\texttt{SensingProcess} typeclass with \texttt{PRESENT}. A sensing
process detects presence: it determines whether two records belong to
the same component. The \texttt{PRESENT} certificate is the union-find
output: two records are present in the same component when their
representatives match. The certificate does not erase the records'
distinct identities; it asserts a relation between them.

The \texttt{GaugeProcess} and \texttt{MEASURABLE} typeclasses extend
this. A gauge process is a sensing process plus a clock plus an event
predicate. The clock provides timestamps; the event predicate decides
which events count. The combination produces \texttt{MEASURABLE}
records: events that are connected in the union-find sense, timestamped
by the clock, and selected by the predicate. \texttt{MEASURABLE} is
the chapter's algorithmic analog of an experimental state.

The connection to the chapter's question is direct. Combining multiple
records is union; finding which component a record belongs to is find.
The two operations are the algorithmic primitives of multi-record
combination. The data structure preserves the records' identities
while exposing their joint structure. The chapter's claim is that
this is the algorithmic form forced by the requirement of admissible
multi-record combination: any algorithm that combines records while
preserving identities reduces to union-find.

\section{Path Compression}
\label{se:path-compression}

After a thousand union operations, the union-find data structure for a
thousand nodes contains internal pointer chains of varying depths. Some
chains are short (the node points directly to the root). Some are
longer (the node points to an intermediate node, which points to
another, which eventually points to the root). The depth determines the
cost of \texttt{find}: a deeper chain means more pointer dereferences.

Path compression is a small modification to \texttt{find} that
flattens the chains as a side effect. The modified procedure walks the
chain from the input node to the root, just like the unmodified
procedure. But on the way back, it updates every node in the chain to
point directly to the root. The next \texttt{find} call on any node in
the same chain will reach the root in one step.

The implementation is a few lines of code:
\begin{verbatim}
def find(x):
    if parent[x] == x:
        return x
    root = find(parent[x])  # recurse to find the root
    parent[x] = root         # compress the path
    return root
\end{verbatim}

The compression is amortized over many operations. The first
\texttt{find} on a deeply nested chain pays the full cost, but it
flattens the chain in the process. Subsequent \texttt{find} calls on
the same nodes are nearly free. With path compression and union-by-
rank combined, the amortized cost of $m$ operations on $n$ nodes is
$O(m \cdot \alpha(n))$, where $\alpha$ is the inverse Ackermann
function. For all practical values of $n$, $\alpha(n) \le 4$. The
amortized cost is effectively constant.

The relevance to the chapter's question is that path compression
preserves the component identity while reducing the cost of
queries. The canonical representative (the root) is the same before
and after compression; only the pointer structure changes. A node $x$
that pointed to $y$, which pointed to $z$ (the root), now points
directly to $z$. The set of nodes in the component is unchanged. The
representative is unchanged. The component's identity is preserved.

This is the chapter's central observation about combination. When
multiple records are merged into a component, the merge preserves
the records' identities. The data structure that represents the merge
(the union-find tree, possibly flattened by path compression) is
internal to the algorithm. Different representations of the same
component produce the same external behavior: the same set of nodes,
the same representative, the same answers to \texttt{find} queries.

The compiler's analog is the namespace caching used during
elaboration. When the elaborator resolves a name like
\texttt{Foo.Bar.baz}, it walks the namespace tree from the root.
Once resolved, the name is cached in a flat lookup table. Subsequent
references to \texttt{Foo.Bar.baz} hit the cache directly,
bypassing the tree walk. The cache plays the same structural role
as path compression: a flatter representation of the same
identity. The name resolves to the same definition; only the
lookup cost changes. (Lean's actual implementation uses scoped
environments and resolution caches rather than a disjoint-set
forest. The point of the comparison is the discipline of
identity-preserving alias merge, not the data structure choice.)

Path compression is also the canonical example of amortized
analysis. The amortized cost of each operation is small, even
though individual operations may be expensive. The cost is amortized
over the lifetime of the data structure: the first $\texttt{find}$
on a long chain is expensive, but it pays for many subsequent
cheap finds. The total cost of $m$ operations is bounded; the cost
per operation is not.

This matters for the chapter's argument because admissibility is
about the structural properties of the algorithm, not the per-
operation cost. The union-find data structure with path compression
is admissible because the component identities are preserved across
all operations. The amortized cost is a quality of the
implementation, not of the algorithmic form.

\section{Compiler Symbol Table}
\label{se:compiler-symbol-table}

Lean's namespace system is a small union-find. Write:
\begin{verbatim}
namespace Foo
  def bar := 3
end Foo
\end{verbatim}
The compiler creates a name \texttt{Foo.bar} in the global namespace.
The name \texttt{bar} alone, used inside the \texttt{namespace Foo}
block, refers to the same definition. The two names are aliases for
the same underlying entity. The compiler maintains this alias
relationship in its symbol table.

When the user opens a namespace:
\begin{verbatim}
open Foo
\end{verbatim}
the compiler adds, to the current scope, mappings from each name in
\texttt{Foo} to its fully qualified form. After this, the user can
write \texttt{bar} alone in the open scope, and the compiler resolves
it to \texttt{Foo.bar}. The opening is the union: the names in
\texttt{Foo} are merged with the names in the current scope, while
preserving their original qualifications.

This is union with preserved source. The merged name \texttt{Foo.bar}
is the same definition it always was. The local alias \texttt{bar} is
a new pointer to that definition; the original name \texttt{Foo.bar}
continues to work. Both names resolve to the same value. The union
does not erase or rename; it adds an additional pointer.

The compiler's analog of \texttt{find} is name resolution. When the
elaborator encounters the bare name \texttt{bar}, it walks the chain
of namespace mappings to find the underlying definition. The chain
is: the current scope's local definitions, then any open
namespaces, then enclosing namespaces, then the global root. The
walk is bounded by the depth of the chain, which is generally small.

The Lean file \texttt{Episode5.lean} defines this in algorithmic
form. The \texttt{Equivalation} typeclass formalizes the equivalence
relation between names: two names are equivalent when they resolve
to the same underlying definition. \texttt{SOURCE} is the original
source name; \texttt{Encoding} is the form used after one or more
namespace operations; \texttt{Abstraction} is the underlying
definition.

When the elaborator sees a use of a name, it computes the
\texttt{Equivalation} class --- the set of names equivalent to the
input. The class is determined by the namespace operations that have
been performed so far. Each \texttt{open} statement is a union; each
name resolution is a find. The data structure that maintains the
namespace is, internally, a union-find with extra annotations for
qualifying prefixes.

The \texttt{CompiledProcess} typeclass marks a definition as
\texttt{EXECUTED}: the full elaboration, type-checking, and
compilation have been performed. An \texttt{EXECUTED} definition can
be used by other definitions; it has reached its final form. The
chain \texttt{SOURCE} $\to$ \texttt{Encoding} $\to$
\texttt{EXECUTED} is the union-find merge of three stages of the
compilation pipeline.

The reason this matters is that the compiler treats the three stages
as distinct records of the same underlying definition. A bug in the
elaborator might produce an \texttt{Encoding} that does not match
the \texttt{SOURCE}. A bug in the typechecker might admit an
\texttt{EXECUTED} value with an incorrect type. The three records
are kept separately so that the compiler can verify their
consistency. The union-find structure provides the framework: each
stage is a node, and the consistency check is whether the three
nodes are in the same component (representing the same definition).

This is also why source maps matter in production compilers. A
runtime error in the compiled output should be traceable back to the
original source line that produced it. The source map is the
\texttt{find} operation that, given an \texttt{EXECUTED} runtime
event, returns the \texttt{SOURCE} location in the original file. The
mapping is union-find: the executed output is one record, the source
is another, they are in the same component, and the mapping retrieves
the source record from the executed record.

The chapter's claim is that this structure is forced by the
chapter's question. The compiler must combine multiple records of
the same definition (source, encoded, executed) without collapsing
them. Union-find is the algorithmic primitive that supports such
combination. Path compression is the optimization that keeps the
queries fast. The combination of the two is the algorithmic
backbone of every compiler that maintains source-to-output mappings.

\section{Horizon as Union Failure}
\label{se:horizon-as-union-failure}

In general relativity, a black hole horizon is a causal boundary. No
signal from inside the horizon can reach an observer outside. The
horizon is one-way: matter and information can fall in but cannot escape.
The interior and the exterior of the horizon are two separate causal
regions, with no admissible exchange of information between them.

This is, in the chapter's algorithmic terms, a union-find failure. The
record of events inside the horizon and the record of events outside
the horizon cannot be unioned. The events on the two sides belong to
different components, and no algorithmic operation can merge them. The
\texttt{find} operation, applied to an interior event and an exterior
event, must return different representatives, because no admissible
chain of inter-event references crosses the horizon.

This is not a defect of the algorithm. It is a structural fact about
the physics. The horizon is real; the disconnect between interior and
exterior is real. The union-find algorithm operating on the causal
record must respect this disconnect by refusing to merge the two
components. A bad algorithm would silently merge them and produce a
unified record where no unification is admissible. The chapter's
algorithm must report the failure.

The same structure appears in less exotic settings. Two computers on
different sides of a network partition cannot communicate. From each
computer's perspective, the other is unreachable. The partition is
the horizon; the two sets of reachable computers are the two
components. Distributed algorithms that operate across the partition
must handle the disconnect: they may not pretend that the partition
does not exist. The CAP theorem (Brewer's theorem) is the formal
statement: a distributed system cannot simultaneously provide
consistency, availability, and partition tolerance. In the presence
of a partition (the horizon), the algorithm must choose to sacrifice
either consistency or availability. The choice is the system's
calibration for the horizon scenario.

Yet another example: in a parallel computation, two threads writing
to different memory locations cannot directly observe each other's
writes without some explicit synchronization. The memory model
imposes a horizon between the threads, and the synchronization
primitives (locks, atomics, memory barriers) are the algorithmic
union operations that selectively cross the horizon. Without the
synchronization, the two threads' records remain in separate
components. With it, the records are merged at the synchronization
point.

The Lean analog in \texttt{Episode4.lean} is the \texttt{Jar} and
\texttt{MeesaProcess} typeclasses. A \texttt{Jar} carries an unresolved
state: a fact that has been recorded but not yet spent (not yet
resolved into a definite value). The \texttt{MeesaProcess} is the
typeclass for processes that handle unresolved states without
prematurely resolving them. The combination produces \texttt{GUNGAN}
records: facts held in suspension, awaiting either resolution or
expiration.

In the horizon analogy, the \texttt{Jar} is the record of an event
that may or may not be in the same component as another event. The
component membership is uncertain --- perhaps a measurement is
required, perhaps the horizon will not be crossed, perhaps the
partition will heal. The \texttt{Jar} carries the uncertainty
admissibly: it does not assert that the event is in some specific
component, but it does carry the event as a candidate for eventual
merging.

The \texttt{MeesaProcess} is the algorithm that handles the
\texttt{Jar}. It may perform a union (if the conditions for merging
are met). It may continue to wait. It may eventually time out and
declare the merger inadmissible. The process's behavior is determined
by its calibration and the structure of the surrounding records. The
chapter's algorithmic claim is that any process handling cross-
component records must instantiate \texttt{MeesaProcess}: it must
admit unresolved states, carry them honestly, and resolve them only
when the underlying conditions warrant.

The chapter's point is that union failures are not exceptions to the
algorithmic framework. They are normal cases that the framework must
handle. A union-find data structure that crashes when given two
records from different components is broken. A union-find data
structure that silently merges everything is also broken. The
correct behavior is to report the component membership honestly:
records in the same component can be merged; records in different
components cannot. The \texttt{find} operation returns the
representative; comparison of representatives reveals component
membership.

The horizon, the network partition, the memory race, and the
\texttt{Jar} are all instances of the same algorithmic structure:
the union-find with explicit component boundaries. The boundaries
are admissible records in their own right. The chapter's
contribution is to make this explicit: a union failure is a
recorded fact, not an exception.

\section{GUNGAN States in Lean}
\label{se:gungan-states-in-lean}

\texttt{Episode4.lean} defines a typeclass cascade for unresolved states.
The cascade starts with \texttt{Jar} (the record of an unresolved fact),
proceeds through \texttt{MeesaProcess} (the process that handles the
\texttt{Jar}), and culminates in \texttt{GUNGAN} (the certificate that
the unresolved state has been carried admissibly).

The \texttt{Jar} structure has the shape:

\begin{leanexcerpt}{Jar}
inductive Jar
  | color         : Fact $\to$ Area $\to$ Jar
  | bang          : Fact $\to$ Jar $\to$ Jar
  | superposition : Fact $\to$ Jar $\to$ Jar $\to$ Jar
\end{leanexcerpt}

A \texttt{Jar} is an inductive with three constructors:
\texttt{color} carries a base \texttt{Fact} and an \texttt{Area};
\texttt{bang} extends a prior \texttt{Jar} under a new \texttt{Fact};
\texttt{superposition} pairs two prior \texttt{Jar}s under a single
\texttt{Fact}. None of the constructors expose a resolved-value flag.
The \texttt{Jar} can be passed between functions while its interior
remains an unresolved layered shape; functions that need a specific
case must pattern-match on the constructors.

The \texttt{MeesaProcess} typeclass:

\begin{leanexcerpt}{MeesaProcess}
structure MeesaProcess
    (Value : Type)
    (Carrier : CarrierProcess Value)
    [DISTINGUISHABLE Value Carrier] ... [MEASURABLE Value Carrier]
  where
  gauge_process : GaugeProcess Value Carrier
  concept       : Jar
  life_debt?    : Jar $\to$ Jar := ...
\end{leanexcerpt}

\texttt{MeesaProcess} is a structure carrying the full project cascade
through \texttt{MEASURABLE}, a \texttt{gauge\_process}, a current
\texttt{concept} (the \texttt{Jar} being held), and a function
\texttt{life\_debt?} that rewrites a \texttt{Jar} under the process.
The function is the algorithmic decision rule; the cascade is the
proof obligation that the rule sits on every prior class admissibly.
The structure formalizes the responsibility of handling \texttt{Jar}s
without premature consumption.

The \texttt{GUNGAN} typeclass:

\begin{leanexcerpt}{GUNGAN}
class GUNGAN
    (Value : Type)
    (Carrier : CarrierProcess Value)
    [DISTINGUISHABLE Value Carrier] ... [MEASURABLE Value Carrier]
  where
  meesa_process : MeesaProcess Value Carrier
  correllant?   : Jar $\to$ Jar $\to$ Prop := ...
\end{leanexcerpt}

\texttt{GUNGAN} certifies that a \texttt{MeesaProcess} handles
\texttt{Jars} admissibly: every unresolved \texttt{Jar} is carried
without modification; every resolution is preceded by a check that the
resolution conditions are met. The certificate is a proof obligation
that the compiler enforces at elaboration time.

The chapter's algorithmic claim is that the combination of \texttt{Jar},
\texttt{MeesaProcess}, and \texttt{GUNGAN} is the typeclass formalization
of admissible handling of cross-component records. A record that may or
may not belong in a given component is a \texttt{Jar}. The process that
manages it is a \texttt{MeesaProcess}. The certificate that the
management is admissible is \texttt{GUNGAN}.

This is the union-find algorithm with explicit handling of uncertain
component membership. A union operation can be conditional on a
\texttt{may\_resolve} check: only resolve the \texttt{Jar} into the
component when the conditions are met. The check may be a calibration
event (a measurement that resolves the uncertainty), or a timeout (the
record expires unresolved), or a manual intervention (the operator
decides). The typeclass formalizes all three.

The compiler enforces the typeclass at elaboration. If a user writes
a function that consumes a \texttt{Jar} without checking
\texttt{resolved}, the compiler refuses to admit the function unless
the function has a \texttt{GUNGAN} certificate justifying the
shortcut. The constraint propagates through the program: every
function that handles \texttt{Jars} must either be \texttt{GUNGAN}-
certified or must defer the handling to a function that is.

The Lean machinery is a working instance of the algorithmic
framework. The typeclasses are not decorative; they are the
algorithm's commitment to handling unresolved states correctly. The
compiler's elaboration is the verification that the commitment is
kept. Programs that violate the commitment do not compile. Programs
that respect it are certified to handle \texttt{Jars} admissibly.

The chapter's closing observation is that this same discipline
applies to physical experiments. An experiment that produces
unresolved data (e.g., a measurement below the instrument's
sensitivity threshold) must report the unresolved status. The data
analysis that follows must respect the unresolved marker. Pretending
the data is resolved produces overconfident conclusions. Honest
analysis carries the unresolved data forward, with appropriate
uncertainty propagation. The Lean \texttt{Jar} and \texttt{GUNGAN}
machinery is the formal version of this experimental discipline.

\begin{bridge}

The chapter's question was how multiple records combine without
collapsing into one undifferentiated record.

The answer is union-find. The algorithm maintains, for each record, a
pointer to a representative of its component. Two records are in the
same component when their representatives match. The records'
identities are preserved; the components are the equivalence classes.
Path compression flattens the pointer chains and makes the queries
fast. The compiler's symbol table is a working instance: source,
encoding, and execution are merged into the same component while
preserving the stage-specific records. Horizons (causal boundaries,
network partitions, memory races) are union failures: components
that cannot be merged. The \texttt{Jar}/\texttt{MeesaProcess}/
\texttt{GUNGAN} typeclass cascade in Lean is the formalization of
admissible handling of unresolved component membership.

What remains is the question of distance. Records that have been
combined into a joint structure are not yet equipped with a notion of
how close or far apart they are. The next chapter develops distance
as a count of refinements --- the number of bisection steps required
to locate one record relative to another. The chapter's algorithmic
forms are bisection and the spline.

\end{bridge}

\begin{coda}{Merging Transit Passes}

The Regional Transit Authority occupies a modern building in the
downtown core. Inside, three agencies that previously operated
independently --- the City Bus Authority, the Suburban Light Rail
District, and the Inter-Regional Express Commission --- are being
combined into a single unified system. The merger has been politically
negotiated for two years. The technical implementation is now
beginning.

Each of the three predecessor agencies has its own fare card. The
city bus card is a plastic card with a magnetic stripe; it carries a
balance in dollars and a usage history. The light rail card is a
contactless smart card; it carries a balance in fare units and a
zone-by-zone history. The inter-regional card is a paper ticket
with a QR code; it carries a route number, a date, and a single-use
flag. Three different physical media, three different account
structures, three different histories.

The unified pass must merge all three.

The project manager is a transit engineer with twenty-three years of
experience. She has been hired specifically for this merger. She has
worked through similar consolidations in two other regions; she
knows where the technical pitfalls are. Her first principle is the
chapter's first principle: the merger must preserve source identity.
The city bus card's history must remain auditable as city bus
history. The light rail card's zones must remain mappable to the
specific light rail stations they encoded. The inter-regional
ticket's route number must remain a valid route number. None of the
three predecessor records may be dissolved into a generic transit
record.

This is the union without collapse. The new unified pass is a
component in the union-find sense: it groups three predecessor
identities into one component, while keeping each identity recoverable.
The unified pass's identifier (a new sixteen-digit number) is the
representative. The predecessor identifiers (the city bus card number,
the light rail card number, the inter-regional ticket number) are
preserved as members of the component.

The technical work begins with the data migration. Each predecessor
agency has a database of cards. The city bus has 1.2 million cards.
The light rail has 800,000. The inter-regional has 200,000 active
tickets at any time, with several million issued historically. The
total is roughly 2.2 million unique passenger records.

The merger script walks the three databases in parallel. For each
predecessor card, it generates a unified pass identifier and records
the mapping: unified ID 0001 corresponds to city bus card
\#XYZ123, light rail card \#ABC456, and inter-regional tickets
issued to user 78901. The mapping table grows by one row per
predecessor card. After the migration completes, the mapping table
has approximately 4 million rows: one per predecessor identity, all
mapping to about 2.2 million unified passes.

The mapping is the union-find data structure. Each row is a node;
the unified pass identifier is the root. Path compression flattens
the lookups: a query for ``which unified pass corresponds to city
bus card \#XYZ123?'' returns the unified pass identifier in
constant time. A query for ``which predecessor records belong to
unified pass 0001?'' returns the three predecessor identifiers,
again in constant time using the reverse index.

The historical records are not migrated. They are linked. The city
bus's database of three years of past trips remains in the city
bus database. The light rail's database remains. The inter-
regional ticket database remains. The unified system queries the
predecessor databases through the mapping table when historical
data is needed. The historical records keep their source identity;
they are not duplicated into the unified system.

This is the SOURCE / EXECUTED distinction at scale. The predecessor
records are the SOURCE; the unified pass is the EXECUTED record.
The Equivalation between them is the mapping table. The compiler
analogy is direct: just as the Lean elaborator maintains the
mapping from source code to elaborated terms, the transit system
maintains the mapping from predecessor cards to unified passes.

After the migration completes, the unified pass goes live. Cardholders
receive their new physical cards in the mail. The new cards are
contactless smart cards with embedded chips. The chip contains the
unified pass identifier. Tap the card on a reader; the reader looks
up the unified pass identifier in the system; the system returns the
balance, the eligibility, and the recent usage history.

The system handles three types of fare calculation. For a city bus
trip, the unified system queries the predecessor city bus fare
schedule and deducts the appropriate amount. For a light rail trip,
it queries the predecessor light rail fare schedule. For an inter-
regional trip, it queries the predecessor inter-regional schedule.
The schedules remain separate; the queries are routed through the
union-find structure to the correct predecessor.

A new feature is cross-modal transfers: a passenger boards a city
bus, transfers to light rail, and continues on the inter-regional
express. The unified system tracks this as a single trip across
three modes. The fare for the multi-modal trip is computed by a new
algorithm that respects the predecessor schedules' interactions: the
combined fare is the maximum of the modal fares, not the sum (a
deliberate policy to encourage multi-modal travel). The new
algorithm is part of the unified system; it does not modify the
predecessor schedules.

Six months after launch, an issue arises. A small subset of users ---
those who had two city bus accounts (one personal, one corporate
expense reimbursement) --- have two unified passes. The migration
script created separate components for the two city bus accounts,
not realizing they were the same person. The user receives two
physical cards, two monthly statements, two balance inquiries.

The fix is a merge operation. The unified system implements a
\texttt{union} call: given two unified pass identifiers, merge them
into one component. The merged pass inherits the union of the
historical mappings. After the merge, both predecessor city bus
accounts point to the same unified pass. The user receives a single
card going forward.

The merge does not erase the prior records. Both city bus accounts
remain in the system, both mapped to the same unified pass. The
unified pass now has two city bus predecessors in its component
membership. The audit trail records both accounts and their merger
date. If a billing dispute arises later about which corporate
expenses belong to which account, the audit trail can disentangle
them.

The chapter's structures are visible throughout. The unified pass
is the union. The predecessor cards are the members. The mapping
table is the union-find data structure. The path compression
ensures fast lookups. The merge operation handles the GUNGAN case:
the two accounts were always the same person, but the system did
not know it until the user reported it. The MeesaProcess held the
inconsistency; the GUNGAN certificate, in this case, was the
manual report from the user. The merge resolved the state.

Some predecessor records cannot be unioned with anything. An
inter-regional ticket that expired before the merger never gets a
unified pass; it was a single-use record that the predecessor
agency had already retired. The expired tickets remain in the
inter-regional database, accessible through legacy queries but not
mapped to the unified system. They are the horizon failures: records
that exist but cannot be merged into the current system.

The project manager closes the migration phase at the eighteen-
month mark. The unified pass is operational. The audit trail is
complete. The component structure is sound: three predecessor
agencies merged into one component per passenger, with predecessor
identities preserved and queryable. The Regional Transit Authority
now operates as a single system, but the source histories remain.

\end{coda}
